\documentclass[11pt,a4paper]{article}

% Document layout and typography.
\usepackage[utf8]{inputenc}
\usepackage[T1]{fontenc}
\usepackage[english]{babel}
\usepackage{lmodern}
\usepackage[a4paper,margin=2.2cm]{geometry}
\usepackage{amsmath,amssymb}
\usepackage{booktabs}
\usepackage{longtable}
\usepackage{array}
\usepackage{xcolor}
\usepackage{hyperref}
\usepackage{listings}

\setlength{\parindent}{0pt}
\setlength{\parskip}{0.55em}

\definecolor{codegray}{RGB}{245,245,245}
\definecolor{darkblue}{RGB}{25,65,120}

\hypersetup{
    colorlinks=true,
    linkcolor=darkblue,
    urlcolor=darkblue,
    citecolor=darkblue,
    pdftitle={FibDisp -- Theory and User Guide},
    pdfauthor={Davide Facciala}
}

\lstset{
    basicstyle=\ttfamily\small,
    backgroundcolor=\color{codegray},
    frame=single,
    breaklines=true,
    columns=fullflexible,
    keepspaces=true
}

\newcommand{\FFT}{\mathrm{FFT}}
\newcommand{\GDD}{\mathrm{GDD}}
\newcommand{\TOD}{\mathrm{TOD}}
\newcommand{\FWHM}{\mathrm{FWHM}}
\newcommand{\SPM}{\mathrm{SPM}}
\newcommand{\TL}{\mathrm{TL}}
\newcommand{\dd}{\mathrm{d}}
\newcommand{\sech}{\operatorname{sech}}

\title{\textbf{FibDisp}\\
\large Simulation of Ultrashort-Pulse Propagation in Gas-Filled Hollow-Core Fibers\\
\large Theory, Numerical Model, Installation, and User Guide}
\author{Davide Faccial\`a}
\date{}

\begin{document}

\maketitle
\tableofcontents
\newpage

% ============================================================================
\section{Purpose and Scope}
% ============================================================================

\texttt{FibDisp} simulates the propagation of an ultrashort laser pulse in a
single-mode gas-filled hollow-core fiber or capillary. The propagation is
calculated with a symmetric split-step Fourier method applied to a generalized
nonlinear Schr\"odinger equation containing:

\begin{itemize}
    \item second-order dispersion, described by $\beta_2$;
    \item third-order dispersion, described by $\beta_3$;
    \item linear attenuation, described by $\alpha$;
    \item Kerr self-phase modulation, described by the nonlinear coefficient
    $\gamma$;
    \item self-steepening, also called the optical-shock correction;
    \item input quadratic and cubic spectral phase, specified by GDD and TOD.
\end{itemize}

The program also provides:

\begin{itemize}
    \item Gaussian, hyperbolic-secant, and super-Gaussian input pulses;
    \item calculated gas/waveguide propagation coefficients and a Manual mode;
    \item selectable propagation terms for diagnostic calculations;
    \item two numerical solvers for the self-steepening nonlinear step;
    \item temporal and spectral evolution maps along the fiber;
    \item an optional fixed post-fiber GDD display in the main Results tab;
    \item spectral-phase polynomial fitting with spectral-power weighting;
    \item pure-GDD compressor analysis and numerical GDD optimization;
    \item parameter sweeps over physical parameters and direct propagation
    coefficients;
    \item compressed and uncompressed pulse-power maps during parameter sweeps;
    \item snapshots of the settings used for the current and most recent sweep;
    \item numerical export in NPZ and CSV formats.
\end{itemize}

The graphical interface contains five tabs:

\begin{enumerate}
    \item \textbf{Settings};
    \item \textbf{Start \& Results};
    \item \textbf{Phase Analysis};
    \item \textbf{GDD Compressor Design};
    \item \textbf{Parameter Sweep}.
\end{enumerate}

The model is intended for studying Kerr-driven spectral broadening, linear
dispersion, attenuation, self-steepening, spectral phase, and the degree to
which an externally applied quadratic spectral phase can compress the pulse.

% ============================================================================
\section{Notation, Coordinates, and Units}
% ============================================================================

The principal quantities used throughout this document are defined here before
they enter the propagation equations.

\begin{longtable}{p{0.18\textwidth}p{0.72\textwidth}}
\toprule
\textbf{Symbol} & \textbf{Definition}\\
\midrule
$z$ & Propagation coordinate along the fiber, in metres.\\
$L$ & Total fiber length, in metres. The GUI displays it in centimetres.\\
$t$ & Laboratory time, in seconds.\\
$T$ & Retarded time, $T=t-\beta_1 z$, in seconds.\\
$U(z,T)$ & Dimensionless complex pulse envelope used by the numerical model.\\
$A(z,T)$ & Power-normalized complex envelope, $A=\sqrt{P_0}\,U$, with units
$\sqrt{\mathrm{W}}$.\\
$P(z,T)$ & Instantaneous pulse power, $P=P_0|U|^2$, in watts.\\
$P_0$ & Peak-power normalization associated with the transform-limited input
shape, in watts.\\
$E_{\mathrm{p}}$ & Pulse energy, $E_{\mathrm{p}}=\int P(T)\,\dd T$, in joules.\\
$\lambda_0$ & Carrier wavelength, in metres; displayed in nanometres.\\
$\nu_0=f_0$ & Carrier ordinary frequency, $\nu_0=c/\lambda_0$, in hertz.\\
$\omega_0$ & Carrier angular frequency, $\omega_0=2\pi\nu_0$, in rad/s.\\
$f_{\FFT}$ & Frequency coordinate generated by the centred NumPy FFT, in hertz.\\
$\nu$ & Physical optical frequency, $\nu=\nu_0-f_{\FFT}$, in hertz.\\
$\omega$ & Physical optical angular frequency, $\omega=2\pi\nu$, in rad/s.\\
$\Omega$ & Physical angular-frequency detuning, $\Omega=\omega-\omega_0=-2\pi f_{\FFT}$.\\
$\widetilde U$ & Fourier-domain envelope corresponding to $U(T)$.\\
$\beta(\omega)$ & Modal propagation constant, in rad/m.\\
$\beta_j$ & $j$th derivative of $\beta$ with respect to $\omega$ at $\omega_0$:
$\beta_j=(\partial^j\beta/\partial\omega^j)_{\omega_0}$.\\
$\beta_1$ & Inverse group velocity, in s/m.\\
$\beta_2$ & Group-velocity dispersion coefficient, in s$^2$/m; displayed as fs$^2$/m.\\
$\beta_3$ & Third-order dispersion coefficient, in s$^3$/m; displayed as fs$^3$/m.\\
$\alpha$ & Power attenuation coefficient, in m$^{-1}$. A field amplitude therefore
contains the factor $\exp(-\alpha z/2)$ when $\alpha$ is constant.\\
$\gamma$ & Kerr nonlinear coefficient, in W$^{-1}$m$^{-1}$.\\
$R$ & Capillary radius, in metres; displayed in micrometres.\\
$A_{\mathrm{eff}}$ & Effective modal area, in m$^2$.\\
$n_{\mathrm{gas}}$ & Linear refractive index of the filling gas.\\
$n_2$ & Nonlinear Kerr refractive index, in m$^2$/W for a specified pressure.\\
$p$ & Gas-pressure value entered in the GUI. Built-in coefficients use this value
in the pressure-scaling relations of the gas model.\\
\bottomrule
\end{longtable}

The speed of light in vacuum is

\begin{equation}
 c=299792458\ \mathrm{m/s}.
\end{equation}

\subsection{Fourier convention}

The centred numerical transform corresponds to the continuous convention

\begin{equation}
\widetilde U(f_{\FFT})
=
\int_{-\infty}^{+\infty}
U(T)e^{-i2\pi f_{\FFT}T}\,\dd T.
\end{equation}

Consequently,

\begin{equation}
\frac{\partial}{\partial T}
\longleftrightarrow
i2\pi f_{\FFT},
\end{equation}

while the physical optical detuning used in plots and phase formulas is

\begin{equation}
\boxed{
\Omega=-2\pi f_{\FFT}.
}
\end{equation}

The physical optical-frequency axis shown by the GUI is

\begin{equation}
\boxed{
\nu=\nu_0-f_{\FFT}.
}
\end{equation}

When the spectral axis is displayed as wavelength,

\begin{equation}
\lambda=\frac{c}{\nu},
\end{equation}

and only points with positive physical frequency are used.

% ============================================================================
\section{Physical Background}
% ============================================================================

\subsection{Kerr refractive index and self-phase modulation}

For an instantaneous Kerr response, the refractive index is written

\begin{equation}
 n=n_0+n_2 I,
\end{equation}

where $n_0$ is the linear refractive index, $n_2$ is the Kerr nonlinear index,
and $I$ is optical intensity in W/m$^2$. For a guided mode with effective area
$A_{\mathrm{eff}}$, intensity and power are related approximately by

\begin{equation}
 I(T)\simeq\frac{P(T)}{A_{\mathrm{eff}}}.
\end{equation}

The nonlinear propagation constant can be expressed with the nonlinear
coefficient

\begin{equation}
\boxed{
\gamma=\frac{n_2\omega_0}{cA_{\mathrm{eff}}}.
}
\end{equation}

If depletion and loss are neglected over a short distance $z$, the Kerr phase
is approximately

\begin{equation}
\phi_{\mathrm{NL}}(T,z)\simeq\gamma P(T)z.
\end{equation}

The instantaneous optical-frequency shift generated by this temporal phase is

\begin{equation}
\boxed{
\Delta\omega(T)
=-\frac{\partial\phi_{\mathrm{NL}}}{\partial T}.
}
\end{equation}

The same relation in ordinary frequency is

\begin{equation}
\boxed{
\Delta\nu(T)
=-\frac{1}{2\pi}
\frac{\partial\phi_{\mathrm{NL}}}{\partial T}.
}
\end{equation}

The leading and trailing edges of a smooth pulse therefore acquire opposite
frequency shifts, producing self-phase-modulation spectral broadening.

A commonly used measure of accumulated nonlinear phase is the $B$-integral,

\begin{equation}
 B=\int_0^L\gamma P_{\mathrm{peak}}(z)\,\dd z,
\end{equation}

where $P_{\mathrm{peak}}(z)$ is the maximum instantaneous power at position $z$.
The $B$-integral is dimensionless and is expressed in radians of nonlinear phase.

\subsection{Linear dispersion}

The modal propagation constant is expanded around $\omega_0$ as

\begin{equation}
\beta(\omega_0+\Omega)
=
\beta_0+\beta_1\Omega
+\frac{\beta_2}{2}\Omega^2
+\frac{\beta_3}{6}\Omega^3+\cdots,
\end{equation}

where

\begin{equation}
\beta_0=\beta(\omega_0),
\qquad
\beta_1=\left.\frac{\partial\beta}{\partial\omega}\right|_{\omega_0},
\qquad
\beta_2=\left.\frac{\partial^2\beta}{\partial\omega^2}\right|_{\omega_0},
\qquad
\beta_3=\left.\frac{\partial^3\beta}{\partial\omega^3}\right|_{\omega_0}.
\end{equation}

The carrier phase associated with $\beta_0$ is not part of the envelope dynamics.
The first-order term is removed by using retarded time $T=t-\beta_1z$.
FibDisp therefore propagates the residual second- and third-order dispersion
through the scalar coefficients $\beta_2$ and $\beta_3$ evaluated at the carrier.

\subsection{Self-steepening}

For a sufficiently broadband pulse, the frequency dependence of the nonlinear
coupling cannot be treated as exactly constant. To first order,

\begin{equation}
\gamma(\omega)
\simeq
\gamma_0\left(1+\frac{\Omega}{\omega_0}\right),
\end{equation}

where $\gamma_0=\gamma(\omega_0)$. The nonlinear spectral term is therefore

\begin{equation}
 i\gamma_0
 \left(1+\frac{\Omega}{\omega_0}\right)
 \widetilde{|A|^2A}.
\end{equation}

Using the Fourier convention defined above, this becomes in retarded time

\begin{equation}
\boxed{
 i\gamma_0|A|^2A
 -\frac{\gamma_0}{\omega_0}
 \frac{\partial}{\partial T}\left(|A|^2A\right).
}
\end{equation}

The first term is SPM. The derivative term is self-steepening. For positive
$n_2$, the nonlinear group-delay correction is larger in the high-power part
of the pulse, producing an intensity-dependent temporal delay and an optical
shock tendency on the trailing side of the pulse.

\subsection{Why spectral broadening can enable compression}

Let the fiber output spectrum be

\begin{equation}
\widetilde U_{\mathrm{out}}(\Omega)
=
\left|\widetilde U_{\mathrm{out}}(\Omega)\right|
 e^{i\phi_{\mathrm{out}}(\Omega)},
\end{equation}

where $\phi_{\mathrm{out}}(\Omega)$ is the output spectral phase. Spectral
broadening increases the bandwidth and can therefore support a shorter
transform-limited pulse. The actual output is usually longer because
$\phi_{\mathrm{out}}$ is not flat.

An external phase-only compressor with transfer function

\begin{equation}
H(\Omega)=e^{i\phi_{\mathrm{comp}}(\Omega)}
\end{equation}

produces

\begin{equation}
U_{\mathrm{comp}}(T)
=
\mathcal F^{-1}
\left\{
\widetilde U_{\mathrm{out}}(\Omega)H(\Omega)
\right\}.
\end{equation}

A pure-GDD compressor uses

\begin{equation}
\boxed{
\phi_{\mathrm{comp}}(\Omega)=\frac{G}{2}\Omega^2,
}
\end{equation}

where $G$ is the applied GDD in fs$^2$ when $\Omega$ is expressed in rad/fs.
A value $G$ that has the opposite sign to the dominant quadratic part of the
output phase can reduce the pulse duration.

% ============================================================================
\section{Input Pulse Model}
% ============================================================================

The user specifies the transform-limited intensity FWHM
$\tau_{\FWHM}$, the pulse energy $E_{\mathrm p}$, and one of three temporal
shapes. A shape-dependent characteristic time $T_0$ is calculated so that the
requested $\tau_{\FWHM}$ is reproduced.

\subsection{Gaussian pulse}

The transform-limited power is

\begin{equation}
P_{\TL}(T)=P_0\exp\left[-\left(\frac{T}{T_0}\right)^2\right],
\end{equation}

and the field envelope is

\begin{equation}
U_{\TL}(T)=\exp\left[-\frac{T^2}{2T_0^2}\right].
\end{equation}

FibDisp uses

\begin{equation}
\boxed{
T_0=\frac{\tau_{\FWHM}}{1.665},
}
\end{equation}

where $1.665$ is the numerical FWHM factor used by the code. The energy
normalization is

\begin{equation}
\boxed{
P_0=\frac{E_{\mathrm p}}{\sqrt{\pi}T_0}.
}
\end{equation}

\subsection{Hyperbolic-secant pulse}

The transform-limited power is

\begin{equation}
P_{\TL}(T)=P_0\sech^2\left(\frac{T}{T_0}\right),
\end{equation}

with field

\begin{equation}
U_{\TL}(T)=\sech\left(\frac{T}{T_0}\right).
\end{equation}

The intensity FWHM gives

\begin{equation}
\boxed{
T_0=
\frac{\tau_{\FWHM}}
{2\operatorname{acosh}(\sqrt{2})},
}
\end{equation}

and, because $\int_{-\infty}^{+\infty}\sech^2(T/T_0)\,\dd T=2T_0$,

\begin{equation}
\boxed{
P_0=\frac{E_{\mathrm p}}{2T_0}.
}
\end{equation}

\subsection{Super-Gaussian pulse}

The transform-limited power is

\begin{equation}
P_{\TL}(T)
=P_0\exp\left[-\left(\frac{T}{T_0}\right)^4\right],
\end{equation}

with field

\begin{equation}
U_{\TL}(T)
=\exp\left[-\frac{1}{2}\left(\frac{T}{T_0}\right)^4\right].
\end{equation}

The characteristic time is

\begin{equation}
\boxed{
T_0=\frac{\tau_{\FWHM}}{2(\ln2)^{1/4}}.
}
\end{equation}

Using

\begin{equation}
\int_{-\infty}^{+\infty}
\exp[-(T/T_0)^4]\,\dd T
=2T_0\Gamma(5/4),
\end{equation}

where $\Gamma$ is the gamma function, the peak-power normalization is

\begin{equation}
\boxed{
P_0=\frac{E_{\mathrm p}}{2T_0\Gamma(5/4)}.
}
\end{equation}

\subsection{Input GDD and TOD}

After constructing the transform-limited field, FibDisp applies the requested
input spectral phase

\begin{equation}
\boxed{
\phi_{\mathrm{in}}(\Omega)
=
\frac{\GDD_{\mathrm{in}}}{2}\Omega^2
+
\frac{\TOD_{\mathrm{in}}}{6}\Omega^3.
}
\end{equation}

Here $\GDD_{\mathrm{in}}$ is in fs$^2$, $\TOD_{\mathrm{in}}$ is in fs$^3$,
and $\Omega$ is expressed in rad/fs for the numerical phase evaluation. The
chirped input spectrum is

\begin{equation}
\widetilde U_{\mathrm{in}}(\Omega)
=
\widetilde U_{\TL}(\Omega)
\exp[i\phi_{\mathrm{in}}(\Omega)].
\end{equation}

Because this transfer function has unit magnitude, input GDD and TOD change the
phase and temporal shape without changing the pulse energy or spectral power.

% ============================================================================
\section{Gas and Hollow-Capillary Model}
% ============================================================================

\subsection{Built-in gases}

The built-in gas choices are Helium, Neon, Argon, Krypton, Xenon, and Nitrogen.
For these gases the program calculates $\beta_2$, $\beta_3$, a
frequency-dependent capillary loss $\alpha(\nu)$, and the nonlinear coefficient
$\gamma$ from the selected gas, pressure, wavelength, and capillary radius.

The effective modal area is approximated by

\begin{equation}
\boxed{
A_{\mathrm{eff}}=0.4766\,\pi R^2.
}
\end{equation}

\subsection{Gas refractive index and pressure scaling}

For each gas a wavelength-dependent reference refractive-index model is used.
The pressure dependence is applied through a Lorentz--Lorenz form. If
$n_{\mathrm{ref}}(\lambda)$ is the reference index and

\begin{equation}
\chi(\lambda,p)
=
\frac{n_{\mathrm{ref}}^2(\lambda)-1}
     {n_{\mathrm{ref}}^2(\lambda)+2}\,p,
\end{equation}

then the gas index used by the code is

\begin{equation}
\boxed{
n_{\mathrm{gas}}(\lambda,p)
=
\sqrt{\frac{1+2\chi}{1-\chi}}.
}
\end{equation}

The pressure variable $p$ is the numerical pressure value entered in the GUI.
The nonlinear index is also scaled linearly with this pressure value in the
built-in model.

\subsection{Large-core capillary approximation}

The fundamental capillary-mode transverse eigenvalue is

\begin{equation}
 u_{11}\simeq2.405.
\end{equation}

The exact ideal-capillary propagation constant can be written

\begin{equation}
\beta(\omega)
=
\sqrt{
 n_{\mathrm{gas}}^2(\omega)\frac{\omega^2}{c^2}
 -\frac{u_{11}^2}{R^2}
}.
\end{equation}

FibDisp evaluates dispersion through the corresponding large-core expansion of
the effective index,

\begin{equation}
\boxed{
n_{\mathrm{eff}}(\lambda)
\simeq
n_{\mathrm{gas}}(\lambda)
\left[
1-
\frac{1}{2}
\left(
\frac{u_{11}\lambda}
{2\pi R n_{\mathrm{gas}}(\lambda)}
\right)^2
\right].
}
\end{equation}

The wavelength derivatives of $n_{\mathrm{eff}}$ are evaluated numerically in
a local wavelength interval around $\lambda_0$. The resulting carrier
coefficients are

\begin{equation}
\boxed{
\beta_2
=
\frac{\lambda_0^3}{2\pi c^2}
\left.
\frac{\dd^2 n_{\mathrm{eff}}}{\dd\lambda^2}
\right|_{\lambda_0},
}
\end{equation}

and

\begin{equation}
\boxed{
\beta_3
=
-\frac{\lambda_0^4}{4\pi^2c^3}
\left[
3\frac{\dd^2 n_{\mathrm{eff}}}{\dd\lambda^2}
+
\lambda_0\frac{\dd^3 n_{\mathrm{eff}}}{\dd\lambda^3}
\right]_{\lambda_0}.
}
\end{equation}

Thus the displayed $\beta_2$ and $\beta_3$ already contain both gas and
waveguide contributions. During propagation they are used as scalar coefficients
in the Taylor expansion around $\omega_0$; the full $\beta(\omega)$ is not
propagated directly.

\subsection{Nonlinear coefficient}

For the built-in gases the nonlinear coefficient is

\begin{equation}
\boxed{
\gamma
=
\frac{n_{2,1}\,p\,\omega_0}
{cA_{\mathrm{eff}}},
}
\end{equation}

where $n_{2,1}$ is the gas nonlinear-index coefficient stored by the model per
unit pressure and $p$ is the pressure value entered in the GUI. Increasing
pressure therefore increases the Kerr coupling approximately linearly while it
also changes the gas contribution to dispersion.

\subsection{Capillary attenuation}

For built-in gases the attenuation is frequency dependent. Let

\begin{equation}
q=\frac{n_v}{n_{\mathrm{gas}}(\lambda_0,p)},
\end{equation}

where $n_v=1.45356$ is the wall-index value used by the capillary-loss model.
At physical optical frequency $\nu$, the attenuation used in the spectral
operator is

\begin{equation}
\boxed{
\alpha(\nu)
=
\left(\frac{u_{11}c}{2\pi\nu}\right)^2
\frac{1}{R^3}
\frac{q^2+1}{\sqrt{q^2-1}}.
}
\end{equation}

The Settings tab displays $\alpha(\nu_0)$, the value at the carrier, while the
propagation uses the frequency-dependent array $\alpha(\nu)$.

\subsection{Manual propagation coefficients}

Manual mode bypasses the built-in coefficient calculation. The user directly
specifies

\begin{equation}
\beta_2\ [\mathrm{fs^2/m}],\qquad
\beta_3\ [\mathrm{fs^3/m}],\qquad
\alpha\ [\mathrm{m^{-1}}],\qquad
\gamma\ [\mathrm{W^{-1}m^{-1}}].
\end{equation}

In Manual mode $\alpha$ is constant across the simulated spectrum. The direct
coefficients are otherwise used in the same propagation equation as the
built-in gas coefficients.

% ============================================================================
\section{Propagation Equation}
% ============================================================================

\subsection{Linear spectral operator}

With retarded time $T=t-\beta_1z$ and physical detuning
$\Omega=\omega-\omega_0$, the linear part of the implemented model is

\begin{equation}
\boxed{
\left.
\frac{\partial\widetilde U}{\partial z}
\right|_{\mathrm L}
=
D(\Omega)\widetilde U,
}
\end{equation}

with

\begin{equation}
\boxed{
D(\Omega)
=
-\frac{\alpha(\omega_0+\Omega)}{2}
+i\frac{\beta_2}{2}\Omega^2
+i\frac{\beta_3}{6}\Omega^3.
}
\end{equation}

For Manual mode, $\alpha$ is constant. For a built-in gas,
$\alpha(\omega_0+\Omega)$ is the capillary-loss function described above.

\subsection{Nonlinear envelope equation}

The numerical envelope is normalized according to

\begin{equation}
A(z,T)=\sqrt{P_0}\,U(z,T),
\qquad
P(z,T)=P_0|U(z,T)|^2.
\end{equation}

The nonlinear part of the propagation equation is

\begin{equation}
\boxed{
\left.
\frac{\partial U}{\partial z}
\right|_{\mathrm{NL}}
=
 i\gamma P_0|U|^2U
-
\frac{\gamma P_0}{\omega_0}
\frac{\partial}{\partial T}\left(|U|^2U\right).
}
\end{equation}

Combining linear and nonlinear contributions gives the exact form used by the
split-step implementation:

\begin{equation}
\boxed{
\frac{\partial U}{\partial z}
=
\mathcal F^{-1}
\left\{
D(\Omega)\widetilde U
\right\}
+
 i\gamma P_0|U|^2U
-
\frac{\gamma P_0}{\omega_0}
\frac{\partial}{\partial T}\left(|U|^2U\right).
}
\label{eq:gnlse}
\end{equation}

If $\alpha$ is constant, Eq.~\eqref{eq:gnlse} can equivalently be written in
time-domain differential form as

\begin{align}
\frac{\partial U}{\partial z}
={}&
-\frac{\alpha}{2}U
-i\frac{\beta_2}{2}
\frac{\partial^2U}{\partial T^2}
+\frac{\beta_3}{6}
\frac{\partial^3U}{\partial T^3}
\nonumber\\
&+
 i\gamma P_0|U|^2U
-
\frac{\gamma P_0}{\omega_0}
\frac{\partial}{\partial T}\left(|U|^2U\right).
\end{align}

The physical roles are:

\begin{center}
\begin{tabular}{p{0.34\textwidth}p{0.52\textwidth}}
\toprule
\textbf{Contribution} & \textbf{Physical effect}\\
\midrule
$-\alpha U/2$ & Linear field attenuation when $\alpha$ is constant.\\[0.4em]
$-i(\beta_2/2)\partial^2U/\partial T^2$ & Group-velocity dispersion.\\[0.4em]
$(\beta_3/6)\partial^3U/\partial T^3$ & Third-order dispersion.\\[0.4em]
$i\gamma P_0|U|^2U$ & Kerr self-phase modulation.\\[0.4em]
$-(\gamma P_0/\omega_0)\partial_T(|U|^2U)$ & Self-steepening / optical-shock term.\\
\bottomrule
\end{tabular}
\end{center}

Each contribution can be enabled or disabled in the Start \& Results tab.
The self-steepening term is physically part of the broadband Kerr response;
enabling it without SPM is mainly useful as a numerical diagnostic.

% ============================================================================
\section{Numerical Self-Steepening Solvers}
% ============================================================================

FibDisp offers two numerical treatments of the nonlinear substep when
self-steepening is enabled.

\subsection{Fast exponential solver}

Define

\begin{equation}
Q(U)=\frac{\partial}{\partial T}\left(|U|^2U\right).
\end{equation}

The nonlinear equation can formally be grouped as

\begin{equation}
\left.
\frac{\partial U}{\partial z}
\right|_{\mathrm{NL}}
=
\left[
 i\gamma P_0|U|^2
-
\frac{\gamma P_0}{\omega_0}
\frac{Q(U)}{U}
\right]U.
\end{equation}

The fast solver freezes the nonlinear coefficients over one global SSFM step
$\Delta z$ and applies multiplicative exponentials. The shock ratio is
regularized as

\begin{equation}
\frac{Q(U)}{U+\epsilon},
\qquad
\epsilon=0.002,
\end{equation}

where $U$ is dimensionless and has a characteristic peak amplitude of order
unity before strong reshaping. The temporal derivative is evaluated
spectrally. This method is computationally efficient and is suitable when the
nonlinear evolution is adequately resolved by the global propagation step.

\subsection{Accurate adaptive RK4 solver}

The accurate solver integrates the nonlinear right-hand side directly,

\begin{equation}
F(U)
=
 i\gamma P_0|U|^2U
-
\frac{\gamma P_0}{\omega_0}
\frac{\partial}{\partial T}\left(|U|^2U\right),
\end{equation}

with classical fourth-order Runge--Kutta stages. The derivative
$\partial_T(|U|^2U)$ is recomputed at every RK4 stage, so no division by
$U+\epsilon$ is required.

Each global nonlinear interval $\Delta z$ is automatically divided into
$N_{\mathrm{sub}}$ internal steps. Let

\begin{equation}
I_U=\max_T |U(T)|^2,
\qquad
g=|\gamma P_0|I_U.
\end{equation}

The substep count is selected so that the estimated nonlinear phase increment
per internal step does not exceed approximately $0.20$ rad and the estimated
self-steepening Courant number does not exceed approximately $0.25$:

\begin{equation}
\frac{g\Delta z}{N_{\mathrm{sub}}}\lesssim0.20,
\end{equation}

\begin{equation}
\frac{3g\Delta z}
{\omega_0\Delta T\,N_{\mathrm{sub}}}
\lesssim0.25,
\end{equation}

where $\Delta T$ is the temporal grid spacing. The factor $3g/\omega_0$ is an
estimate of the intensity-characteristic speed associated with the reduced
self-steepening equation and is used only as a numerical step-size criterion.

The accurate mode is slower but is substantially more robust when the nonlinear
coefficient, pulse peak power, or optical-shock strength is large. If more than
5000 internal RK4 steps would be required inside one global SSFM interval,
FibDisp stops the calculation and asks for a finer propagation or temporal grid.

No time integrator can compensate for an unresolved temporal shock or a
spectrum reaching the FFT boundaries. Strong self-steepening calculations
should therefore also be checked for convergence with respect to $N_Z$, FFT
size, and temporal window.

% ============================================================================
\section{Split-Step Fourier Propagation}
% ============================================================================

The full propagation uses a symmetric split-step Fourier method. Let
$\widehat D$ denote the linear propagation operator and $\widehat N$ the
nonlinear operator. Over one interval $\Delta z$,

\begin{equation}
U(z+\Delta z)
\simeq
 e^{\widehat D\Delta z/2}
 e^{\widehat N\Delta z}
 e^{\widehat D\Delta z/2}
 U(z).
\end{equation}

The numerical sequence is:

\begin{enumerate}
    \item Fourier transform the temporal field;
    \item multiply the spectrum by $\exp[D(\Omega)\Delta z/2]$;
    \item inverse transform to retarded time;
    \item apply the nonlinear substep with the selected self-steepening solver;
    \item Fourier transform again;
    \item apply the second factor $\exp[D(\Omega)\Delta z/2]$;
    \item inverse transform to obtain the field at the next propagation point.
\end{enumerate}

If $N_Z$ is the number of points on the propagation axis, then there are
$N_Z-1$ intervals and

\begin{equation}
\boxed{
\Delta z=\frac{L}{N_Z-1}.
}
\end{equation}

The first point is $z=0$ and the final point is $z=L$.

% ============================================================================
\section{Temporal and Spectral Numerical Grids}
% ============================================================================

The FFT contains

\begin{equation}
\boxed{N=2^n}
\end{equation}

samples, where $n$ is the GUI value \texttt{Points time/freq 2\^{}n}.

Let $m$ denote the GUI value \texttt{ratio time/FWHM m}. In the implemented
grid construction, $m$ multiplies the shape-dependent characteristic time
$T_0$. The periodic temporal interval is

\begin{equation}
\boxed{
-mT_0\le T<mT_0,
}
\end{equation}

with total duration

\begin{equation}
T_{\mathrm{win}}=2mT_0.
\end{equation}

The temporal spacing is

\begin{equation}
\boxed{
\Delta T=\frac{2mT_0}{N},
}
\end{equation}

and the FFT-frequency spacing is

\begin{equation}
\boxed{
\Delta f=\frac{1}{N\Delta T}
=\frac{1}{T_{\mathrm{win}}}.
}
\end{equation}

At fixed $m$, increasing $n$ increases the temporal resolution and the
available Nyquist bandwidth. At fixed $n$, increasing $m$ enlarges the temporal
window but also increases $\Delta T$; strong broadening can therefore require
increasing both $m$ and $n$.

The \texttt{Saved z samples} value controls only how many propagation snapshots
are retained for the temporal/spectral evolution map. It does not change
$N_Z$, $\Delta z$, or the propagation physics.

% ============================================================================
\section{Computational-Boundary Warnings}
% ============================================================================

During propagation, FibDisp monitors the first and last samples of the temporal
and spectral fields. A warning is generated when an edge amplitude exceeds
$10^{-3}$ of the corresponding initial peak field amplitude. In symbols, a
frequency-boundary warning is triggered when

\begin{equation}
\frac{|\widetilde U_{\mathrm{edge}}|}
     {\max|\widetilde U_{\mathrm{in}}|}>10^{-3},
\end{equation}

and the analogous test is applied in time.

A boundary warning indicates that periodic FFT wrap-around may influence the
solution. The usual remedies are:

\begin{itemize}
    \item increase $n$ when the spectrum approaches the frequency boundaries;
    \item increase $m$ when the temporal pulse approaches the time boundaries;
    \item increase both when a larger temporal window is needed without losing
    temporal resolution.
\end{itemize}

Parameter sweeps perform the same check for every sweep point and report the
boundary type and approximate propagation position at which the condition was
first detected.

% ============================================================================
\section{Derived Quantities}
% ============================================================================

\subsection{Output energy and transmission}

The calculated spectral power density is denoted $S(\nu)$ and has units W/Hz.
The output energy is

\begin{equation}
\boxed{
E_{\mathrm{out}}
=\int S_{\mathrm{out}}(\nu)\,\dd\nu.
}
\end{equation}

The transmission percentage is

\begin{equation}
\boxed{
\eta
=100\frac{E_{\mathrm{out}}}{E_{\mathrm{in}}},
}
\end{equation}

where $E_{\mathrm{in}}$ is the input pulse energy.

\subsection{Spectral broadening factor}

For a spectral density $S(f)$, define the power-weighted mean frequency

\begin{equation}
\bar f
=
\frac{\int fS(f)\,\dd f}{\int S(f)\,\dd f}
\end{equation}

and standard deviation

\begin{equation}
\sigma_f
=
\sqrt{
\frac{\int(f-\bar f)^2S(f)\,\dd f}{\int S(f)\,\dd f}
}.
\end{equation}

FibDisp reports the dimensionless broadening factor

\begin{equation}
\boxed{
B=\frac{\sigma_{f,\mathrm{out}}}{\sigma_{f,\mathrm{in}}}.
}
\end{equation}

The sign reversal between $f_{\FFT}$ and physical detuning does not change this
standard-deviation ratio.

\subsection{Transform-limited output pulse}

The transform-limited field associated with the output spectral amplitude is

\begin{equation}
\boxed{
U_{\TL}(T)
=\mathcal F^{-1}
\left\{
|\widetilde U_{\mathrm{out}}(\Omega)|
\right\}.
}
\end{equation}

Its temporal FWHM is the duration obtainable if the complete output spectral
phase could be removed while preserving the spectral amplitude.

\subsection{Temporal instantaneous-frequency shift}

Let

\begin{equation}
\phi_T(T)=\arg U(T)
\end{equation}

be the unwrapped temporal envelope phase. The displayed chirp is the ordinary
frequency shift

\begin{equation}
\boxed{
\Delta\nu(T)
=-\frac{1}{2\pi}
\frac{\dd\phi_T}{\dd T}.
}
\end{equation}

FibDisp stores this quantity in hertz and displays it in terahertz.

\subsection{Spectral group delay}

Let $\phi(\omega)$ be the unwrapped spectral phase. The group delay is

\begin{equation}
\boxed{
GD(\omega)=\frac{\dd\phi}{\dd\omega}.
}
\end{equation}

It is stored in seconds and displayed in femtoseconds. With the internal FFT
coordinate, the same derivative is evaluated as

\begin{equation}
GD
=-\frac{1}{2\pi}
\frac{\dd\phi}{\dd f_{\FFT}}.
\end{equation}

% ============================================================================
\section{Tab 1: Settings}
% ============================================================================

The Settings tab defines the baseline parameters used for a propagation or a
parameter sweep.

\subsection{Gas and fiber parameters}

\begin{longtable}{p{0.29\textwidth}p{0.62\textwidth}}
\toprule
\textbf{GUI field} & \textbf{Meaning}\\
\midrule
Select gas & Built-in gas model or Manual direct-coefficient mode.\\
Pressure (atm) & Pressure value $p$ used by the built-in gas model.\\
Fiber length (cm) & Total propagation distance $L$.\\
Radius ($\mu$m) & Capillary radius $R$.\\
\bottomrule
\end{longtable}

\subsection{Pulse parameters}

\begin{longtable}{p{0.29\textwidth}p{0.62\textwidth}}
\toprule
\textbf{GUI field} & \textbf{Meaning}\\
\midrule
Pulse shape & Gaussian, Sech, or Super-Gaussian transform-limited intensity profile.\\
Energy (mJ) & Input pulse energy $E_{\mathrm p}$.\\
TL FWHM (fs) & Transform-limited intensity FWHM $\tau_{\FWHM}$.\\
Wavelength (nm) & Carrier wavelength $\lambda_0$.\\
GDD (fs$^2$) & Input quadratic spectral-phase coefficient $\GDD_{\mathrm{in}}$.\\
TOD (fs$^3$) & Input cubic spectral-phase coefficient $\TOD_{\mathrm{in}}$.\\
\bottomrule
\end{longtable}

\subsection{Simulation parameters}

\begin{longtable}{p{0.29\textwidth}p{0.62\textwidth}}
\toprule
\textbf{GUI field} & \textbf{Meaning}\\
\midrule
Points prop. axis z & Number $N_Z$ of propagation-axis samples.\\
Points time/freq $2^n$ & FFT sample count $N=2^n$.\\
ratio time/FWHM $m$ & Numerical half-window factor; the implemented half-width is $mT_0$.\\
\bottomrule
\end{longtable}

\subsection{Propagation coefficients}

For a built-in gas, the tab displays $\beta_2$, $\beta_3$, $\alpha(\nu_0)$,
and $\gamma$. The values update when gas, pressure, wavelength, or radius is
changed. The coefficient fields are read-only in this mode.

In Manual mode the four fields become editable. The value entered for $\alpha$
is then used as a frequency-independent loss coefficient.

\subsection{Confirm (preview input)}

The preview button generates the input pulse using the selected shape, energy,
FWHM, wavelength, GDD, TOD, and numerical grid. It displays temporal power and
phase together with spectral power density and phase before fiber propagation.

% ============================================================================
\section{Tab 2: Start \& Results}
% ============================================================================

This tab performs the propagation and displays the principal output diagnostics.

\subsection{Scalar results}

The reported scalar values are:

\begin{itemize}
    \item output pulse energy $E_{\mathrm{out}}$;
    \item transmission $\eta$;
    \item spectral broadening factor $B$;
    \item transform-limited output FWHM.
\end{itemize}

\subsection{Selectable propagation terms}

The following terms can be independently enabled or disabled:

\begin{itemize}
    \item SPM;
    \item self-steepening;
    \item GVD ($\beta_2$);
    \item TOD ($\beta_3$);
    \item loss ($\alpha$).
\end{itemize}

The self-steepening solver can be selected as either \textbf{Fast
(exponential)} or \textbf{Accurate RK4 (adaptive)}. The solver choice has no
effect when self-steepening is disabled.

\subsection{Raw and fixed-GDD result display}

The \textbf{Result display} control has two modes:

\begin{description}
    \item[Raw output] displays the complex field directly at the fiber exit.
    \item[Fixed GDD] applies a user-selected external GDD $G$ to the final
    output spectrum and reconstructs the corresponding temporal field.
\end{description}

The fixed GDD is phase-only post-processing:

\begin{equation}
\widetilde U_G(\Omega)
=
\widetilde U_{\mathrm{out}}(\Omega)
\exp\left(i\frac{G}{2}\Omega^2\right).
\end{equation}

It does not modify the fiber propagation, the saved $z$ evolution, the raw
simulation result used by Phase Analysis, or the GDD Compressor Design tab.

For a fixed-GDD temporal display, the automatic x-range is obtained from the
first and last samples above 1\% of the compensated pulse peak, with additional
padding equal to half of that significant span on each side when the numerical
window permits. This rule affects only the displayed axis limits; it does not
modify the temporal field or the computational time grid.

\subsection{Result plots}

The $2\times3$ layout contains:

\begin{enumerate}
    \item pulse power and temporal phase;
    \item spectral power density and spectral phase;
    \item transform-limited pulse power;
    \item instantaneous-frequency shift (chirp);
    \item spectral group delay;
    \item propagation evolution map.
\end{enumerate}

The spectral axis can be displayed in THz or nm. The Matplotlib navigation
toolbar provides zoom, pan, home, and save functions and remains outside the
resizable plot area so that it stays accessible when the window height is
reduced.

\subsection{Spectral or temporal evolution along the fiber}

The lower-right panel can display either

\begin{equation}
S(\nu,z)
\end{equation}

or

\begin{equation}
P(T,z),
\end{equation}

where $S$ is spectral power density and $P$ is temporal power. Both maps use the
same saved propagation positions. The number of stored positions is controlled
by \texttt{Saved z samples}.

Both maps always represent the field \emph{inside the fiber before any external
fixed-GDD compensation}. An external compressor is applied only after the final
propagation position and therefore has no physical meaning at intermediate $z$.

\subsection{Saving a single simulation}

A simulation can be saved as NPZ or CSV. The NPZ archive retains the numerical
arrays, including complex temporal and spectral fields. CSV export writes the
main table and companion files for the propagation maps.

The saved quantities include:

\begin{itemize}
    \item temporal grid and FFT-frequency grid;
    \item input and output complex fields;
    \item input and output temporal power;
    \item input and output spectral power density;
    \item temporal and spectral phase;
    \item instantaneous-frequency shift;
    \item spectral group delay;
    \item transform-limited pulse;
    \item spectral evolution $S(\nu,z)$;
    \item temporal evolution $P(T,z)$;
    \item physical and numerical settings associated with the run.
\end{itemize}

% ============================================================================
\section{Tab 3: Phase Analysis}
% ============================================================================

The Phase Analysis tab analyzes the \textbf{raw fiber output}. The optional
fixed GDD selected only for the Start \& Results display is not included in this
analysis.

\subsection{Polynomial convention}

The fit variable is physical angular-frequency detuning in rad/fs. The fitted
phase is written

\begin{equation}
\boxed{
\phi_{\mathrm{fit}}(\Omega)
=
p_0
+p_1\Omega
+\frac{p_2}{2!}\Omega^2
+\frac{p_3}{3!}\Omega^3
+\cdots.
}
\end{equation}

The coefficients therefore have direct dispersion units:

\begin{align}
p_0 &\quad [\mathrm{rad}],\\
p_1 &\quad [\mathrm{fs}],\\
p_2 &\quad [\mathrm{fs^2}],\\
p_3 &\quad [\mathrm{fs^3}].
\end{align}

In particular, $p_1$ is a fitted group-delay offset, $p_2$ is the fitted GDD,
and $p_3$ is the fitted TOD.

\subsection{Spectral-power-weighted least squares}

For spectral samples indexed by $i$, let $S_i$ be the output spectral power
density, $\phi_i$ the unwrapped output phase, and $\Omega_i$ the physical
detuning. The fit minimizes

\begin{equation}
\boxed{
\chi^2
=
\sum_i S_i
\left[
\phi_i-\phi_{\mathrm{fit}}(\Omega_i)
\right]^2.
}
\end{equation}

The weights are normalized internally, so only their relative values matter.
Strong spectral regions therefore dominate over weak wings with poorly
conditioned phase.

\subsection{Fit range and display axes}

The fit interval can be entered in THz or nm. Changing the displayed unit
converts the limits so that the same physical spectral interval is preserved.
The fitted quantity is always computed using $\Omega$ in rad/fs.

The spectral-phase plot automatically chooses its y-limits from the phase data
that lie inside the currently displayed x-range. When a polynomial fit is
present, the visible part of the fitted curve is included in the same y-range
calculation. A small margin is added above and below. This display scaling does
not alter phase unwrapping, fit limits, weights, or fitted coefficients.

The fitted GDD is a global, spectral-power-weighted quadratic component over the
selected interval. After strong nonlinear broadening it need not equal the
local curvature

\begin{equation}
\left.
\frac{\dd^2\phi}{\dd\omega^2}
\right|_{\omega_0}.
\end{equation}

% ============================================================================
\section{Tab 4: GDD Compressor Design}
% ============================================================================

This tab applies a pure quadratic spectral phase to the raw fiber output and
reconstructs the temporal field.

\subsection{Pure-GDD transfer function}

For applied GDD $G$,

\begin{equation}
\boxed{
H_G(\Omega)
=\exp\left(i\frac{G}{2}\Omega^2\right),
}
\end{equation}

and

\begin{equation}
\widetilde U_G(\Omega)
=
\widetilde U_{\mathrm{out}}(\Omega)H_G(\Omega).
\end{equation}

Since $|H_G|=1$, the spectral power density is unchanged. Only spectral phase
and the reconstructed temporal field change.

\subsection{Manual GDD}

The user can enter any value of $G$ in fs$^2$ and inspect the compensated pulse,
the raw fiber output, the transform-limited pulse, the spectrum, and the
compensated spectral phase.

\subsection{Temporal quality metrics}

Four metrics are available for characterizing or optimizing the compressed
pulse.

\subsubsection{Main FWHM}

The main FWHM is the temporal distance between the nearest half-maximum
crossings on either side of the global pulse maximum. It describes the main
peak but can ignore separated satellites.

\subsubsection{RMS duration}

For temporal power $P(T)$, define the power-weighted centroid

\begin{equation}
\bar T
=
\frac{\int T P(T)\,\dd T}
{\int P(T)\,\dd T}.
\end{equation}

The RMS duration is

\begin{equation}
\boxed{
\sigma_T
=
\sqrt{
\frac{\int(T-\bar T)^2P(T)\,\dd T}
{\int P(T)\,\dd T}
}.
}
\end{equation}

This metric is sensitive to the complete temporal profile, including weak
pedestals and distant satellites.

\subsubsection{95\% energy width}

The 95\% energy width is the shortest contiguous interval $[T_1,T_2]$ that
contains 95\% of the total temporal energy:

\begin{equation}
\boxed{
\Delta T_{95}=T_2-T_1,
}
\end{equation}

with

\begin{equation}
\frac{\int_{T_1}^{T_2}P(T)\,\dd T}
     {\int_{-\infty}^{+\infty}P(T)\,\dd T}
\ge0.95,
\end{equation}

and $T_2-T_1$ chosen as small as possible.

\subsubsection{95\%-bounded FWHM}

This is the default optimization metric. First, the shortest interval
containing 95\% of the pulse energy is found. Starting from each boundary of
that interval and moving toward the global peak, the first crossing of
$P_{\max}/2$ is located. The distance between those two crossings is the
95\%-bounded FWHM.

The metric remains sensitive to sufficiently strong pre-pulses and post-pulses
inside the 95\%-energy interval while retaining an FWHM-like interpretation.

\subsection{Optimized GDD}

The GDD objective can contain several local minima. FibDisp therefore combines:

\begin{enumerate}
    \item a spectral-phase estimate used to form the seed
    $G\simeq-p_2$;
    \item a broad global GDD scan;
    \item identification of promising minima;
    \item progressively finer local grids around several candidate basins;
    \item retention of the best objective value encountered;
    \item a small final local refinement when useful.
\end{enumerate}

The optimized GDD is the value that minimizes the selected temporal quality
metric. It is an operational compressor setting and is not required to equal
the negative of the fitted GDD when higher-order or structured spectral phase
is important.

% ============================================================================
\section{Tab 5: Parameter Sweep}
% ============================================================================

The Parameter Sweep tab uses the current Settings values, propagation-term
switches, self-steepening solver, and Phase Analysis configuration as the
baseline for a family of calculations.

\subsection{Sweep parameters}

The selectable sweep parameter is one of:

\begin{itemize}
    \item pressure $p$;
    \item fiber length $L$;
    \item capillary radius $R$;
    \item pulse energy $E_{\mathrm p}$;
    \item transform-limited input FWHM $\tau_{\FWHM}$;
    \item carrier wavelength $\lambda_0$;
    \item input GDD;
    \item input TOD;
    \item direct $\beta_2$;
    \item direct $\beta_3$;
    \item direct $\alpha$;
    \item direct $\gamma$;
    \item compensating output GDD.
\end{itemize}

When $\beta_2$, $\beta_3$, $\alpha$, or $\gamma$ is selected as the sweep
parameter, only that coefficient is replaced at each sweep point. All other
baseline physical parameters and direct coefficients remain fixed.

For a built-in gas:

\begin{itemize}
    \item a $\beta_2$, $\beta_3$, or $\gamma$ override leaves the calculated
    frequency-dependent capillary loss $\alpha(\nu)$ unchanged;
    \item an $\alpha$ override replaces the calculated loss with the selected
    constant attenuation value over the whole spectrum.
\end{itemize}

\subsection{Sweep values and spacing}

A sweep is defined by a minimum $x_{\min}$, maximum $x_{\max}$, and number of
values $N_{\mathrm{sweep}}$. Both endpoints are included.

Let

\begin{equation}
u_j=\frac{j}{N_{\mathrm{sweep}}-1},
\qquad
j=0,1,\ldots,N_{\mathrm{sweep}}-1.
\end{equation}

For linear spacing,

\begin{equation}
x(u)=x_{\min}+(x_{\max}-x_{\min})u.
\end{equation}

For exponential spacing with base $b>0$, $b\ne1$,

\begin{equation}
\boxed{
s_{\exp}(u)=\frac{b^u-1}{b-1},
}
\end{equation}

and

\begin{equation}
x(u)=x_{\min}+(x_{\max}-x_{\min})s_{\exp}(u).
\end{equation}

For logarithmic spacing,

\begin{equation}
\boxed{
s_{\log}(u)
=
\frac{\ln[1+(b-1)u]}{\ln b},
}
\end{equation}

with the same mapping to $[x_{\min},x_{\max}]$. Both normalized mappings
approach linear spacing as $b\rightarrow1$ and can be used for signed parameter
ranges because the logarithm or exponential acts on the normalized coordinate
$u$, not directly on the physical parameter.

The Preview field shows the exact list of values before the sweep starts.

\subsection{Compressor criterion}

The sweep optimizer has its own compressor-criterion selector. The available
metrics are:

\begin{itemize}
    \item 95\%-bounded FWHM;
    \item RMS duration;
    \item 95\% energy width;
    \item Main FWHM.
\end{itemize}

The selected criterion is used to determine the optimized GDD and optimized
compressed duration at each relevant sweep point.

\subsection{Quantities calculated at each point}

The sweep result contains, where applicable:

\begin{itemize}
    \item phase-fit GDD;
    \item optimized GDD;
    \item actual raw output FWHM;
    \item transform-limited output duration;
    \item optimized compressed duration;
    \item boundary-warning information;
    \item raw temporal and spectral output fields needed for post-fiber
    compression views.
\end{itemize}

The sweep table reports the optimized compressed duration as a stable numerical
reference. The selectable compressed plots described below may instead display
a fixed or swept GDD without changing the stored raw fiber outputs.

\subsection{Phase-fit range used during a sweep}

If the Phase Analysis tab contains a valid fit range, the sweep uses that
physical range. If no valid range is available, an automatic significant
spectral interval is selected independently for each sweep point. The phase fit
remains spectral-power weighted.

\subsection{Plot choices}

The Plot selector provides:

\begin{itemize}
    \item Phase-fit GDD;
    \item Optimized GDD;
    \item Actual output FWHM;
    \item TL duration;
    \item Compressed duration;
    \item Spectral power map;
    \item Pulse power map (uncompressed);
    \item Pulse power map (compressed).
\end{itemize}

The spectral map displays $S(\nu,x)$ and the temporal maps display $P(T,x)$,
where $x$ denotes the selected sweep parameter. If different sweep points use
different physical frequency or time grids, the maps are interpolated onto a
common grid for visualization. The original per-point grids and arrays are
retained separately in the numerical result.

\subsection{Compression mode for compressed plots}

The \textbf{Compressed plots} selector controls both \textbf{Compressed
duration} and \textbf{Pulse power map (compressed)}.

For every sweep other than a compensating-GDD sweep, the available modes are:

\begin{itemize}
    \item \textbf{Optimized GDD}: use the independently optimized GDD for each
    sweep point;
    \item \textbf{Fixed GDD}: apply the same user-entered GDD to every sweep
    point.
\end{itemize}

If the sweep parameter itself is \textbf{Compensating GDD}, a third mode is
available:

\begin{itemize}
    \item \textbf{Swept GDD}: apply the GDD value represented by each point on
    the sweep axis.
\end{itemize}

Thus a compensating-GDD sweep offers Swept GDD, Optimized GDD, and Fixed GDD;
all other sweeps offer Optimized GDD and Fixed GDD. Changing this selector is a
post-processing operation on the stored complex output spectrum and does not
rerun the fiber propagation.

\subsection{Special propagation strategies}

For a fiber-length sweep, the field is propagated once to the largest requested
length. Whenever a requested intermediate length is reached, the complex field
is analyzed at that physical position. If necessary, the active SSFM interval
is shortened so the requested length is reached exactly.

For a compensating-GDD sweep, the fiber parameters do not change. The fiber is
therefore propagated once, the optimized GDD is determined once for that raw
output, and every swept GDD is then applied to the same complex spectrum.

\subsection{Viewing sweep settings}

Two read-only windows help identify the full parameter context of a sweep:

\begin{description}
    \item[View current settings] shows the live GUI values that would be used
    if a sweep were started at that moment.
    \item[View last run sweep settings] shows a frozen snapshot of the settings
    used by the most recently completed sweep. Editing the Settings tab later
    does not modify this snapshot.
\end{description}

The snapshot includes gas/fiber parameters, pulse parameters, direct
coefficients, propagation-term switches, self-steepening solver, numerical
grid, sweep definition, phase-fit configuration, and compressor settings.

\subsection{Saving sweep data}

Sweep results can be exported as NPZ or CSV. The data include the sweep values,
phase-fit and optimized GDD, durations, boundary warnings, common interpolated
maps, native per-point spectral and temporal axes, raw pulse-power arrays,
complex output spectra for compressor reconstruction, and baseline settings
associated with the sweep.

% ============================================================================
\section{Interpreting Fitted and Optimized GDD}
% ============================================================================

After nonlinear propagation, the output spectral phase can contain quadratic,
cubic, higher-order, and oscillatory contributions:

\begin{equation}
\phi(\Omega)
=
\phi_0+\phi_1\Omega
+\frac{\phi_2}{2}\Omega^2
+\frac{\phi_3}{6}\Omega^3
+\phi_{\mathrm{HO}}(\Omega),
\end{equation}

where $\phi_{\mathrm{HO}}$ denotes contributions beyond the displayed low-order
terms.

Three different quadratic-phase concepts should be distinguished:

\begin{enumerate}
    \item the local spectral-phase curvature
    \[
    \left.\frac{\dd^2\phi}{\dd\omega^2}\right|_{\omega_0};
    \]
    \item the global spectral-power-weighted fitted GDD $p_2$ over a chosen
    spectral interval;
    \item the optimized compressor GDD $G_{\mathrm{opt}}$, defined as the value
    that minimizes a selected temporal quality metric.
\end{enumerate}

For a nearly quadratic phase,

\begin{equation}
G_{\mathrm{opt}}\simeq-p_2.
\end{equation}

After strong broadening, the two can differ because a pure quadratic phase
cannot remove TOD, higher-order spectral phase, or rapid phase structure.

% ============================================================================
\section{Installation and Running from Python}
% ============================================================================

FibDisp consists of two Python source files that must be kept in the same
directory:

\begin{description}
    \item[\texttt{fibdisp\_core.py}] propagation physics, pulse generation,
    phase analysis, compressor calculations, and parameter sweeps;
    \item[\texttt{fibdisp\_gui.py}] Tkinter graphical interface, plotting,
    controls, warnings, and data export.
\end{description}

The principal Python dependencies are NumPy, SciPy, Matplotlib, and Tkinter.
PyInstaller is needed only when creating a standalone executable.

A standalone executable is \emph{not} required to use FibDisp. Running the
Python source directly is often the simplest choice for development and
scientific work.

\subsection{Windows: Python virtual environment}

Install a recent 64-bit Python distribution with Tk support. Python 3.12 is a
suitable choice. Open PowerShell in the directory containing
\texttt{fibdisp\_gui.py} and \texttt{fibdisp\_core.py}.

Create a local virtual environment:

\begin{lstlisting}
python -m venv .venv
\end{lstlisting}

If PowerShell prevents activation of local scripts, allow scripts only for the
current PowerShell process:

\begin{lstlisting}
Set-ExecutionPolicy -Scope Process -ExecutionPolicy Bypass
\end{lstlisting}

Activate the environment:

\begin{lstlisting}
.\.venv\Scripts\Activate.ps1
\end{lstlisting}

Install the numerical packages:

\begin{lstlisting}
python -m pip install --upgrade pip
python -m pip install numpy scipy matplotlib
\end{lstlisting}

Run FibDisp directly:

\begin{lstlisting}
python fibdisp_gui.py
\end{lstlisting}

The packages are installed inside the local \texttt{.venv} directory. They do
not need to be installed system-wide.

\subsection{Linux: Python virtual environment}

On Debian or Ubuntu, install Python, the virtual-environment module, and Tk if
needed:

\begin{lstlisting}
sudo apt update
sudo apt install python3 python3-venv python3-tk
\end{lstlisting}

Create and activate the environment:

\begin{lstlisting}
python3 -m venv .venv
source .venv/bin/activate
\end{lstlisting}

Install dependencies and run:

\begin{lstlisting}
python -m pip install --upgrade pip
python -m pip install numpy scipy matplotlib
python fibdisp_gui.py
\end{lstlisting}

For another Linux distribution, install the equivalent Python and Tk packages
with the system package manager.

\subsection{macOS: Python virtual environment}

Use a Python installation with Tk support, then open Terminal in the FibDisp
source directory:

\begin{lstlisting}
python3 -m venv .venv
source .venv/bin/activate
python -m pip install --upgrade pip
python -m pip install numpy scipy matplotlib
python fibdisp_gui.py
\end{lstlisting}

If several Python installations are present, verify the selected interpreter
with

\begin{lstlisting}
python3 --version
\end{lstlisting}

before creating the environment.

\subsection{Anaconda or Miniconda}

FibDisp can also be run without compiling inside an isolated Conda environment.
From Anaconda Prompt, PowerShell, or a terminal with Conda initialized:

\begin{lstlisting}
conda create -n fibdisp python=3.12 numpy scipy matplotlib tk
conda activate fibdisp
python fibdisp_gui.py
\end{lstlisting}

To use an already existing Conda environment, activate it and install any
missing dependencies:

\begin{lstlisting}
conda activate my_environment
conda install numpy scipy matplotlib tk
python fibdisp_gui.py
\end{lstlisting}

The source directory only needs to contain \texttt{fibdisp\_gui.py} and
\texttt{fibdisp\_core.py}; Python imports the core module from the same folder.

% ============================================================================
\section{Building Standalone Executables}
% ============================================================================

A standalone application bundles the Python interpreter and the required Python
packages. End users can then run FibDisp without installing Python separately.
The recommended PyInstaller format is \texttt{--onedir}, which creates one
application directory containing the executable and its runtime libraries.

Native executables should be built on the target operating system: build the
Windows application on Windows, the Linux application on Linux, and the macOS
application on macOS.

\subsection{Install build dependencies}

Inside a virtual environment or Conda environment containing FibDisp, install
PyInstaller:

\begin{lstlisting}
python -m pip install pyinstaller
\end{lstlisting}

If NumPy, SciPy, and Matplotlib are not already installed, install them as well:

\begin{lstlisting}
python -m pip install numpy scipy matplotlib
\end{lstlisting}

\subsection{Generic PyInstaller command}

Run this command from the directory containing the two FibDisp source files:

\begin{lstlisting}
python -m PyInstaller --noconfirm --clean --windowed --onedir --name FibDisp fibdisp_gui.py
\end{lstlisting}

The \texttt{--windowed} option suppresses a separate terminal window for the
GUI application. The \texttt{--onedir} option keeps all bundled runtime files
inside one application directory.

\subsection{Windows output}

The Windows application is normally created at

\begin{lstlisting}
dist\FibDisp\FibDisp.exe
\end{lstlisting}

Distribute the complete \texttt{dist\textbackslash FibDisp} directory, not only
\texttt{FibDisp.exe}, because the executable uses the bundled libraries stored
beside it.

If a supplied PowerShell build script is used, run it from its own directory:

\begin{lstlisting}
powershell -ExecutionPolicy Bypass -File .\build_windows.ps1
\end{lstlisting}

A local Python virtual environment created for building can be deleted after a
successful build if no further builds are required.

\subsection{Linux output}

The Linux executable is normally created at

\begin{lstlisting}
dist/FibDisp/FibDisp
\end{lstlisting}

Distribute the complete \texttt{dist/FibDisp} directory. For broad Linux
compatibility, build on a distribution at least as old as the oldest
distribution that must run the application, because the resulting binary is
linked to the build system's C runtime.

\subsection{macOS output}

A windowed macOS build normally produces

\begin{lstlisting}
dist/FibDisp.app
\end{lstlisting}

Apple Silicon and Intel binaries should be built and tested for their intended
architectures. An unsigned local application can trigger Gatekeeper. Public
macOS distribution should use appropriate code signing and notarization.

\subsection{Build directories}

PyInstaller commonly creates:

\begin{description}
    \item[\texttt{.venv}] optional local Python environment used for building;
    \item[\texttt{build}] temporary PyInstaller build files;
    \item[\texttt{dist}] final distributable application;
    \item[\texttt{FibDisp.spec}] PyInstaller build specification.
\end{description}

The \texttt{build} directory can be removed after a successful build. The
\texttt{.venv} directory can also be removed if it is not needed for later
rebuilding. The \texttt{dist} output is the application to keep or distribute.

% ============================================================================
\section{Recommended Simulation Workflow}
% ============================================================================

A practical sequence is:

\begin{enumerate}
    \item choose gas, pressure, fiber radius, fiber length, pulse shape, pulse
    energy, transform-limited FWHM, carrier wavelength, and input phase in
    \textbf{Settings};
    \item inspect the displayed $\beta_2$, $\beta_3$, $\alpha$, and $\gamma$;
    \item use \textbf{Confirm (preview input)} to verify the input field and
    numerical window;
    \item select the propagation terms and self-steepening solver in
    \textbf{Start \& Results};
    \item select the number of stored $z$ samples needed for propagation maps;
    \item run the propagation and check boundary warnings;
    \item inspect raw pulse power, spectrum, chirp, group delay, transform
    limit, and the temporal or spectral evolution map;
    \item use the fixed-GDD Result display when a quick post-fiber compensation
    view is useful;
    \item use \textbf{Phase Analysis} to quantify the weighted spectral phase;
    \item use \textbf{GDD Compressor Design} to compare phase-fit GDD with the
    optimized pure-GDD compressor;
    \item choose a temporal optimization metric appropriate to the pulse
    structure, especially when satellites are present;
    \item use \textbf{Parameter Sweep} to study trends and maps while keeping
    track of the full baseline through the settings-view buttons;
    \item repeat selected calculations with finer numerical grids to verify
    convergence before treating small differences as physical.
\end{enumerate}

% ============================================================================
\section{Numerical Convergence Checks}
% ============================================================================

A numerical result should be checked for convergence with respect to the
propagation and FFT grids.

\subsection{Propagation-step convergence}

Increase $N_Z$, thereby reducing $\Delta z$, and compare important observables
such as:

\begin{itemize}
    \item output spectrum and temporal profile;
    \item broadening factor;
    \item transform-limited FWHM;
    \item phase-fit GDD;
    \item optimized GDD;
    \item optimized compressed duration.
\end{itemize}

Strong self-steepening can require substantially finer $z$ resolution in the
fast exponential mode. The accurate RK4 mode reduces this sensitivity by
subdividing the nonlinear step internally, but convergence of the complete
split-step solution should still be checked.

\subsection{FFT-resolution convergence}

Increase $n$ so that $N=2^n$ increases. Verify that spectral and temporal fine
structure, fitted phase, and compressor results do not change significantly.
Frequency-boundary warnings are a direct indication that the available
bandwidth is insufficient.

\subsection{Temporal-window convergence}

Increase $m$ when the pulse, temporal satellites, or dispersive tails approach
the periodic time boundaries. Since increasing $m$ at fixed $N$ increases
$\Delta T$, increase $n$ as well when the temporal resolution must be preserved.

\subsection{Self-steepening convergence}

For calculations dominated by optical-shock dynamics:

\begin{itemize}
    \item compare Fast and Accurate RK4 solvers when practical;
    \item increase $N_Z$;
    \item increase FFT resolution;
    \item verify that the spectrum remains well separated from the FFT edges;
    \item verify that the sharpest temporal feature is sampled by several time
    points rather than one or two grid samples.
\end{itemize}

% ============================================================================
\section{Model Limitations}
% ============================================================================

The implemented propagation model includes Kerr nonlinearity, self-steepening,
loss, and dispersion through $\beta_3$. It does not explicitly include:

\begin{itemize}
    \item photoionization and plasma generation;
    \item ionization loss;
    \item a time-resolved molecular Raman response;
    \item dispersion coefficients beyond $\beta_3$ in the propagation Taylor
    expansion;
    \item propagation with the complete frequency-dependent modal
    $\beta(\omega)$;
    \item multimode transverse dynamics;
    \item longitudinally varying capillary radius or pressure profiles;
    \item real-compressor amplitude response or higher-order compressor phase
    beyond the applied pure quadratic phase.
\end{itemize}

The large-core capillary approximation is used to calculate the modal effective
index for the built-in gas coefficients. Far from the carrier, particularly
when a spectrum spans a very large fraction of an octave, a local
$\beta_2+\beta_3$ Taylor expansion can become less accurate than propagation
based on the full modal dispersion relation.

For regimes where ionization, plasma dynamics, Raman response, multimode
propagation, or higher-order dispersion are important, the result should be
interpreted as the prediction of the model defined by Eq.~\eqref{eq:gnlse}, not
as a complete description of every experimental process.

% ============================================================================
\section{Summary of Sign and Phase Conventions}
% ============================================================================

The principal conventions are:

\begin{itemize}
    \item the centred FFT coordinate is $f_{\FFT}$;
    \item physical optical frequency is
    \[
    \nu=\nu_0-f_{\FFT};
    \]
    \item physical detuning is
    \[
    \Omega=-2\pi f_{\FFT};
    \]
    \item input or compressor quadratic phase is written
    \[
    \phi_{\GDD}(\Omega)=\frac{\GDD}{2}\Omega^2;
    \]
    \item cubic phase is written
    \[
    \phi_{\TOD}(\Omega)=\frac{\TOD}{6}\Omega^3;
    \]
    \item temporal instantaneous-frequency shift is
    \[
    \Delta\nu(T)
    =-\frac{1}{2\pi}\frac{\dd\phi_T}{\dd T};
    \]
    \item spectral group delay is
    \[
    GD=\frac{\dd\phi}{\dd\omega};
    \]
    \item a predominantly positive output quadratic spectral phase is
    compensated by a negative applied GDD, and vice versa.
\end{itemize}

% ============================================================================
\section{Citation}
% ============================================================================

If you use \texttt{FibDisp} in scientific work, please cite \textbf{both the
FibDisp software and the original scientific publication}.

\textbf{Software citation:}

\begin{quote}
D. Faccial\`a, C. Vozzi, G. Sansone, S. De Silvestri, M. Nisoli, and S. Stagira,
\textit{FibDisp}, computer software,\\
\href{https://github.com/omegatype/fibdisp-web}{https://github.com/omegatype/fibdisp-web}.
\end{quote}

The preferred machine-readable software citation is provided in
\texttt{CITATION.cff}.

\textbf{Original scientific publication:}

\begin{quote}
C. Vozzi, M. Nisoli, G. Sansone, S. Stagira, and S. De Silvestri,
``Optimal spectral broadening in hollow-fiber compressor systems,''
\textit{Applied Physics B} \textbf{80}, 285--289 (2005).\\
DOI: \href{https://doi.org/10.1007/s00340-004-1721-1}{10.1007/s00340-004-1721-1}.
\end{quote}

See \texttt{NOTICE.md} for attribution information and
\texttt{CITATION.cff} for the preferred machine-readable software citation
metadata.

% ============================================================================
% ============================================================================
\section{References}
% ============================================================================

\begin{thebibliography}{9}

\bibitem{nisoli1996}
M. Nisoli, S. De Silvestri, and O. Svelto,
``Generation of high energy 10 fs pulses by a new pulse compression
technique,''
\textit{Applied Physics Letters} \textbf{68}, 2793--2795 (1996).
DOI: 10.1063/1.116609.

\bibitem{nisoli1997}
M. Nisoli, S. De Silvestri, O. Svelto, R. Szip\"ocs, K. Ferencz,
C. Spielmann, S. Sartania, and F. Krausz,
``Compression of high-energy laser pulses below 5 fs,''
\textit{Optics Letters} \textbf{22}, 522--524 (1997).

\bibitem{vozzi2005}
C. Vozzi, M. Nisoli, G. Sansone, S. Stagira, and S. De Silvestri,
``Optimal spectral broadening in hollow-fiber compressor systems,''
\textit{Applied Physics B} \textbf{80}, 285--289 (2005).
DOI: 10.1007/s00340-004-1721-1.

\bibitem{schenkel2003}
B. Schenkel, J. Biegert, U. Keller, C. Vozzi, M. Nisoli, G. Sansone,
S. Stagira, S. De Silvestri, and O. Svelto,
``Generation of 3.8-fs pulses from adaptive compression of a cascaded
hollow fiber supercontinuum,''
\textit{Optics Letters} \textbf{28}, 1987--1989 (2003).
DOI: 10.1364/OL.28.001987.

\end{thebibliography}

\end{document}
